\emph{Setting.} The route of Theorem \claimref{synthesis.einstein-equations-from-c-loc-and-clausius}, under Premise (C$_{\mathrm{loc}}$) and its hypotheses (H1) to (H4), shows that on its domain every admissible pair (field, conserved $C^{1}$ source $T$) satisfies its item (4), $\mathcal{R}(a) + 8\pi G_{\mathrm{cl}}T(a) = (\Lambda + 4\pi G_{\mathrm{cl}}\tau)a$ with $\Lambda$ constant and $G_{\mathrm{cl}} = 1/(4\eta)$, the connection being torsion-free (a $C^{3}$ derived metric with its Levi-Civita connection). This is a statement about a solution set. The inverse variational problem asks for a local functional whose Euler-Lagrange equations have that solution set; any invertible operator applied to the equations has the same solution set, so the problem includes the choice of a multiplier. Throughout, $\tr S = \sum_{a} e^{a}\cdot S(e_{a})$ in an orthonormal frame, and in four dimensions $\tr\mathcal{G} = -\mathcal{R}$ (the scalar), $\tr(\Lambda a) = 4\Lambda$.

\emph{The Helmholtz condition and principal symbols.} If $E = \delta S/\delta\phi$ for a local functional $S$, the Fr\'echet derivative $DE$ is the second variation of $S$, a symmetric bilinear form: $\int \ip{k}{DE[h]} = \int \ip{h}{DE[k]}$ for compactly supported $h$, $k$ (the Helmholtz condition; its sufficiency on star-shaped domains is the Volterra-Vainberg theorem, cited from memory and not used here, since every sufficiency claim below has an explicit witness). For a linear operator $P = \sum a^{ij}(x)\partial_{i}\partial_{j} + (\text{lower order})$ on symmetric tensors, the formal adjoint has principal part $\sum (a^{ij})^{\top}\partial_{i}\partial_{j}$, so self-adjointness forces the principal symbol $\sigma(p)$ to be symmetric for the pairing $\ip{k}{h} = k_{ab}h^{ab}$ at every point and every $p$. For a natural second-order expression $E[g]$ the principal symbol of $DE$ at a point depends only on $g$ there; in a normal gauge at the point ($\hb = 1$, $\Omega = 0$ there, Definition \claimref{obj.synthesis.gtg-flat-spacetime-gauge-fields}) it is the flat-background symbol. If $M$ is algebraic in $g$, $D(ME)[h] = M(DE[h]) + (DM[h])E$ and the second term has order zero in $h$, so the principal symbol of $D(ME)$ is $M\circ\sigma(DE)$. A Lorentz-covariant second-order symbol on symmetric $h$ is a combination of $A_{1}, \dots, A_{5}$ (the covariant quadratic forms in $p$ with values in symmetric tensors, linear in $h$: $h$ enters through $\tr h$, $p\cdot h\cdot p$, $hp$ or itself).

\emph{(i)} The check builds the linearised Riemann tensor from $h$ and $p$, $R_{abcd} = \tfrac12(p_{b}p_{c}h_{ad} + p_{a}p_{d}h_{bc} - p_{a}p_{c}h_{bd} - p_{b}p_{d}h_{ac})$ (sphere-positive convention; the overall sign is a common factor and does not affect symmetry, transversality or kernels), contracts it to the Ricci and Einstein symbols and expresses them in the basis, finding $\sigma(\mathcal{G}) = -\tfrac12(A_{1} + A_{2} - A_{3} + A_{4} - A_{5})$ and $\sigma(\Ric) = -\tfrac12(A_{1} + A_{4} - A_{5})$. Symmetry of $\ip{k}{\sigma h}$: $A_{3}$, $A_{4}$, $A_{5}$ are symmetric and $A_{1}$, $A_{2}$ are adjoint to each other ($\ip{k}{A_{1}h} = \tr h\;p\cdot k\cdot p$, $\ip{k}{A_{2}h} = \tr k\;p\cdot h\cdot p$), so the self-adjoint covariant symbols are those with equal $A_{1}$ and $A_{2}$ coefficients, a space of dimension $4$ (computed exactly on the ten basis tensors at two rational momenta). The Einstein symbol has equal coefficients and passes; the Ricci symbol has $-\tfrac12$ and $0$ and fails. Since $\tr\sigma(\mathcal{G})$ is built from $A$-terms whose trace contributes only to $A_{2}$ and $A_{3}$, $M_{\beta}\sigma(\mathcal{G}) = \sigma(\mathcal{G}) + \beta\eta\,\tr\sigma(\mathcal{G})$ changes the $A_{2}$ coefficient by $\beta$ times a nonzero number, so its asymmetry is $\beta$ times that of $M_{1}\sigma(\mathcal{G}) - \sigma(\mathcal{G})$, nonzero; the check confirms the linearity in $\beta$ at $\beta \in \{-\tfrac12, -\tfrac14, \tfrac37\}$, and that $M_{-1/2}\sigma(\mathcal{G}) = \sigma(\Ric)$. Algebraically, $M_{-1/2}(\mathcal{G} + \Lambda + 8\pi G_{\mathrm{cl}}T) = \mathcal{R} - \Lambda + 8\pi G_{\mathrm{cl}}(T - \tfrac12\tau)$, which is the route's item (4) (checked on a sample Ricci tensor and source). Hence $M_{\beta}(\mathcal{G} + \Lambda)$ fails the Helmholtz condition at every point of every field when $\beta \neq 0$. For $\beta = 0$ the Einstein-Hilbert functional with $\Lambda$ has metric Euler-Lagrange expression proportional to $\mathcal{G} + \Lambda$, and the stored functional has frame equation $2(\mathcal{G}(a) + \Lambda a)I$ (Theorem \claimref{synthesis.macdowell-mansouri-principle-fixes-the-trace-free-part}(i)): the condition is met, by explicit witnesses. \emph{Field variable.} All of this is for the metric (or the frame) as the varied field with the pairing $\ip{k}{h}$. If $\beta \neq -\tfrac14$, $M_{\beta}$ is invertible ($1 + 4\beta \neq 0$) and $\psi$ with $\delta g = M_{\beta}(\delta\psi)$ is an equivalent field; the Euler-Lagrange expression of the Einstein-Hilbert functional in $\psi$ is $M_{\beta}^{\top}$ of that in $g$, that is $M_{\beta}(\mathcal{G} + \Lambda)$ up to a factor ($M_{\beta}$ is symmetric for the pairing), and its $\psi$-derivative has principal symbol $M_{\beta}\,\sigma(\mathcal{G})\,M_{\beta}$, symmetric. The check confirms the symmetry at $\beta \in \{-\tfrac12, \tfrac37, -2\}$ and that the Ricci symbol composed with $M_{-1/2}$ on the input is $M_{-1/2}\sigma(\mathcal{G})M_{-1/2}$: the route's Ricci form is the Einstein-Hilbert Euler-Lagrange expression of a densitised metric. So (i) is a statement about the varied field, not about the equations alone.

\emph{(ii)} $M_{-1/4}$ maps $\eta$ to $0$ and has rank $9$ on the ten-dimensional space (checked), so $M_{-1/4}(\mathcal{G} + \Lambda) = \mathcal{G} - \tfrac14\tr(\mathcal{G})$, the trace-free Ricci form, in which $\Lambda$ does not occur; for $\beta \neq -\tfrac14$, $M_{\beta}$ is invertible ($1 + 4\beta \neq 0$) and $M_{\beta}(\mathcal{G} + \Lambda) = 0$ fixes the trace $\mathcal{R}$ in terms of $\Lambda$, so $\Lambda$ is a parameter of the equation. So the trace-free form is the one member of the family with $\Lambda$ free, and by (i) it fails the Helmholtz condition on all variations of the metric. On trace-free $h$, $k$ (variations that fix the volume element; in Gauge Theory Gravity, frame variations at fixed $\det\hb$) the terms $A_{1}$, $A_{2}$, $A_{3}$ pair to zero ($\tr h = \tr k = 0$), so every covariant symbol is self-adjoint there; the check confirms this for the trace-free symbol on a basis of the nine-dimensional trace-free space, and the condition is met there by a witness: the Einstein-Hilbert functional at fixed volume element has $\delta S = \int\ip{\mathcal{G}}{h}$ with $\tr h = 0$, whose Euler-Lagrange expression is the trace-free part of $\mathcal{G}$. With an enlarged field content the trace-free equations are variational on all variations with $\Lambda$ a constant of integration: the Henneaux-Teitelboim form, Fact \claimref{dark-energy.imported-henneaux-teitelboim-promotion}. This is the route's own variational form: on the route $N$ is initial data and its solution space is the promoted one (Theorem \claimref{synthesis.einstein-route-makes-n-initial-data}), and under the promotion the norm mode is one global pair with no local mode (Theorem \claimref{synthesis.promoted-norm-one-global-pair}). For the metric alone a fixed value is needed (see the last paragraph of this step), and the frozen point's constraint carries it as an input, the registered input $N$ fixing $\Lambda$; it does not compute it: with $I = \nInf + \tfrac{\Lambda}{6}\nO$ in $\Cl(2,4)$ ($\nO = \tfrac12(e_{+} + e_{-})$, $\nInf = e_{+} - e_{-}$), $Q(I) = \Lambda/3$ (Theorem \claimref{dark-energy.frozen-point-makes-stelle-west-vector}, checked in $\Cl(2,4)$). For de Sitter curvature at $\Lambda'$ the store's Ricci tensor is $\Lambda' g$ (the convention of Theorem \claimref{synthesis.macdowell-mansouri-principle-fixes-the-trace-free-part}); the trace-free form vanishes for every $\Lambda'$, and $\mathcal{G} + 3Q(I) = (\Lambda - \Lambda')g$ vanishes only at $\Lambda' = \Lambda$ (checked at two values of $\Lambda$ and three of $\Lambda'$). Beyond the family: Lovelock (1971, 1972; Fact \claimref{synthesis.imported-lovelock-and-anderson-duchamp-natural-lagrangians}(a)) proved that in four dimensions a symmetric tensor built naturally from the metric and its derivatives up to second order that is divergence-free (equivalently, for natural expressions, an Euler-Lagrange expression) is $a\mathcal{G} + bg$; its vacuum set has the single value $\Lambda = b/a$, so no natural variational problem of order at most two for the metric alone has the route's solution set with $\Lambda$ free. The order bound is needed: the Weyl-squared functional $\int W\cdot W$ is fourth order (the check finds its second variation at the flat background to be $2\,\Ric\cdot\Ric - \tfrac23 R^{2}$ in the linearised curvature, nonzero and homogeneous of degree four in $p$), and its Euler-Lagrange equations, the Bach equations, hold on every Einstein field of every $\Lambda$ in four dimensions (Fact \claimref{synthesis.imported-conformal-gravity-bach-equations-and-conformal-gauge-form}(a)), so its solution set contains the route's vacuum set with $\Lambda$ free and is larger. It is the frozen-point integrand on the normal tractor connection of a torsion-free field (Theorem \claimref{synthesis.frozen-point-integrand-on-tractor-transport-is-weyl-squared}(ii), (iii)).

\emph{(iii)} Invariance of a functional under displacement gauge transformations (diffeomorphisms) gives, by Noether's second theorem, $\nabla_{a}E^{ab} = 0$ identically; linearised at a flat point this is $p^{a}\sigma(p)_{ab} = 0$ for every $h$. On the route the same identity enters at the Bianchi step of Theorem \claimref{synthesis.einstein-equations-from-c-loc-and-clausius}(4), with $\nabla\cdot T = 0$ from (H3). Solving $p^{a}(\sum_{i} c_{i}A_{i})_{ab} = 0$ exactly gives a plane: $c_{1} = -c_{3}$ and $c_{2} = c_{4} = -c_{5}$; it contains the Einstein symbol and $A_{1} - A_{3} = (p_{a}p_{b} - \eta_{ab}p^{2})\tr h$, which is transverse and not self-adjoint. With self-adjointness ($c_{1} = c_{2}$) one line remains, spanned by the Einstein symbol. This is the linearised form of Lovelock's theorem, proved here without the import.

\emph{(iv)} Let $S$ be as stated. By (i) and (ii) (within the family) and by Lovelock's theorem (in general), its gravitational frame expression is $c(\mathcal{G} + \Lambda)$ with $\Lambda = 3Q(I)$; matching the route's coefficient $8\pi G_{\mathrm{cl}}$ in front of the source normalised as in (v) fixes $c$, that is $G = G_{\mathrm{cl}} = 1/(4\eta)$. If two local functionals have the same Euler-Lagrange expressions, their difference has identically vanishing Euler-Lagrange expression; by the exactness of the variational bicomplex it is locally a divergence, and for natural Lagrangians of a metric in four dimensions the variationally trivial ones are combinations of the Euler and Pontryagin densities plus divergences of natural currents (Anderson and Duchamp 1980, Anderson 1984; Fact \claimref{synthesis.imported-lovelock-and-anderson-duchamp-natural-lagrangians}(b)). Consistently, the check finds that the Euler density $E_{4} = \mathrm{Riem}\cdot\mathrm{Riem} - 4\Ric\cdot\Ric + R^{2}$ has zero second variation at the flat background, and that adding $R^{2}$ to it gives exactly the second variation of $R^{2}$. So in the metric reading $S$ is the Einstein-Hilbert functional with $G_{\mathrm{cl}}$ and $\Lambda$ plus $c_{E}E_{4} + c_{P}P_{4}$ plus a divergence, with $c_{E}$, $c_{P}$ not fixed by the route. The stored frozen-point functional is $(R - 2\Lambda + \alpha E_{4})/16\pi G$ with $\alpha = -3/(4\Lambda)$ (Theorem \claimref{synthesis.euler-term-horizon-entropy-shift}(i)), the member with $c_{E} = \alpha/16\pi G_{\mathrm{cl}}$, $c_{P} = 0$; nothing on the route selects it. In first-order variables, on the field space of connections preserving $I$, the route sees only torsion-free fields, so a density whose Euler-Lagrange expressions vanish whenever the torsion does is invisible to it; the curvature-only classification (Theorem \claimref{synthesis.curvature-only-tractor-actions-are-macdowell-mansouri}(iii), granted for this first-order clause) finds the admissible densities $\kappa\,\epsilon(I, X, F, F) + \kappa' H + \kappa'' K$, with $K$ the Chern-Weil density (no field equations) and $H$ a Holst-type density whose frame equation vanishes on every curvature obeying the first Bianchi identity. Granting that classification, the route fixes $\kappa$ (by $G = G_{\mathrm{cl}}$; the frame equation of $\epsilon(I, X, F, F)$ is the Einstein form, Theorem \claimref{synthesis.macdowell-mansouri-principle-fixes-the-trace-free-part}(i)) and leaves $\kappa'$, $\kappa''$ free; within that class the Euler coefficient is the one carried by $\epsilon(I, X, F, F)$, fixed by $\Lambda$. On the normal tractor connections the same integrand is the Weyl square (Theorem \claimref{synthesis.frozen-point-integrand-on-tractor-transport-is-weyl-squared}(iii)), of fourth order, excluded by the order bound. \emph{For the imports of the open problem:} the route is compatible with a local action at fixed $\Lambda$ (witness: the stored functional) but does not single one out, so the existence of the action is not derived; the coefficients it leaves free, Euler and Pontryagin, are exactly those that distinguish the MacDowell-Mansouri functional from the Einstein-Hilbert functional, so its form is not derived either.

\emph{(v)} Let $\mathcal{S}$ be a local functional of $(g, \psi)$ with Euler-Lagrange system $(c(\mathcal{G} + \Lambda + 8\pi G_{\mathrm{cl}}T), E_{\psi})$, $c$ constant, and put $\mathcal{S}' = \mathcal{S} - S_{m}$. Its $\psi$-expression is $E_{\psi} - E_{\psi} = 0$ identically. The Hessian of $\mathcal{S}'$ is symmetric (Helmholtz), and its $(\psi, g)$ block is $D_{g}$ of the $\psi$-expression, which is $0$; so its $(g, \psi)$ block, $D_{\psi}$ of the $g$-expression, is $0$: $D_{\psi}\big(c(\mathcal{G} + \Lambda + 8\pi G_{\mathrm{cl}}T) - \delta S_{m}/\delta g\big) = 0$. With $\delta S_{m}/\delta g = 8\pi G_{\mathrm{cl}}\,c\,T_{H}$ (the normalisation of $T_{H}$ fixed by the same constant $c$) and $D_{\psi}(\mathcal{G} + \Lambda) = 0$, this is $D_{\psi}(T - T_{H}) = 0$ at every field; on a field space connected in $\psi$, $T - T_{H} = \Theta[g]$ does not depend on the matter, and $\Theta = 0$ if both vanish at $\psi = 0$. Conversely, if $T = T_{H}$ the system is the Euler-Lagrange system of $c$ times the Einstein-Hilbert functional plus $S_{m}$. The check verifies the cross condition $\ip{k}{K(p)} = L(-p)[k]$ for a scalar field with Lagrangian $\sqrt{g}(\tfrac12 g^{ab}\partial_{a}\psi\partial_{b}\psi - V(\psi)) - \tfrac12\xi\psi^{2}\sqrt{g}\,\mathcal{R}^{\circ}$ at the flat background, $K$ the $\psi$-derivative of the metric expression and $L$ the metric derivative of the matter equation, on the ten basis tensors at two momenta: with a background gradient $u$ and $V' = v_{1}$ constant ($\xi = 0$), $\ip{k}{K} = \tfrac12\tr k\,(u\cdot p - v_{1}) - u\cdot k\cdot p = L(-p)[k]$; with a constant background $s$, the $\xi$ terms $\xi s(p^{2}\tr k - p\cdot k\cdot p)$ appear on both sides (from $\sqrt{g}(-f\mathcal{G}^{\circ} - g\Box f + \nabla\nabla f)$, $f = -\tfrac12\xi\psi^{2}$, and from $-\xi\psi\,\delta\mathcal{R}^{\circ}$ with $\delta\mathcal{R}^{\circ} = p\cdot k\cdot p - p^{2}\tr k$). It fails when the source carries an altered trace term $-\tfrac12\lambda\,\tr k\,(u\cdot p)$, $\lambda = 1$, and when the source is the minimal stress while the matter equation carries $-\xi\psi\mathcal{R}^{\circ}$. So the condition ties the heat tensor to the matter functional, minimal or not; whether the matter couples minimally is not decided by it. So the energy part of the matter coupling is equivalent to the identification $T = T_{H}$, not derived from the route. The connection equation (spin density and torsion) does not occur on a torsion-free route, so the spin part of Fact \claimref{synthesis.imported-einstein-cartan-matter-coupling} is not reached.

\emph{Hypotheses that do the work.} Locality and naturality of the functional; the multiplier algebraic in $g$ (a multiplier depending on curvature or derivatives changes the principal symbol and is not covered); Euler-Lagrange expressions of order at most two (for (ii) beyond the family and for (iv)); the frozen point's $\Lambda$ (for the whole field; without it, unimodular variations); for (v), that the matter equations come from a local functional. The route's own premise (C$_{\mathrm{loc}}$) and (H1) to (H4) enter only through its output. Imported: Lovelock 1971 and 1972, and Anderson and Duchamp 1980 and Anderson 1984 on natural variationally trivial Lagrangians (Fact \claimref{synthesis.imported-lovelock-and-anderson-duchamp-natural-lagrangians}); Bach 1921 and the Bach-flatness of Einstein fields in four dimensions (Fact \claimref{synthesis.imported-conformal-gravity-bach-equations-and-conformal-gauge-form}). Literature cited from memory: the sufficiency of the Helmholtz conditions (Volterra-Vainberg), not used; Noether's second theorem.
